\documentclass[12pt, twoside]{article}
\input{../preamble}

\fancyhead[LO]{La Scuola d'Italia / Dr.\ Huson / IB Math \\* 28 May 2026}
\fancyhead[RO]{Markscheme}
\fancyfoot[C]{\thepage}

\begin{document}
\subsubsection*{8.4 Classwork: Rational Functions Challenge --- Markscheme}

\medskip

%================================================================
%--- Problem 1 [8 marks]
%================================================================
\noindent\textbf{1.} $\displaystyle f(x)=\dfrac{3x-1}{x-2}$, \ $x\ne 2$.

\medskip

\textbf{(a)} Vertical asymptote: $x=2$. \hfill \textbf{A1}

\medskip

\textbf{(b)} Horizontal asymptote: $y=3$
\ \ (ratio of leading coefficients $=\tfrac{3}{1}$). \hfill \textbf{A1}

\medskip

\textbf{(c)} $x$-intercept: $3x-1=0$ \ \ $\Rightarrow$ \ \ $x=\dfrac{1}{3}$
\hfill \textbf{A1}

$y$-intercept: $f(0)=\dfrac{-1}{-2}=\dfrac{1}{2}$ \hfill \textbf{A1}

\emph{Accept the ordered-pair forms $\left(\tfrac{1}{3},0\right)$ and
$\left(0,\tfrac{1}{2}\right)$.}

\medskip

\textbf{(d)} Let $y=\dfrac{3x-1}{x-2}$ and swap $x$ and $y$: \hfill \textbf{M1}

$x(y-2)=3y-1$

$xy-3y=2x-1$, \ \ \ $y(x-3)=2x-1$ \hfill \textbf{A1}

$f^{-1}(x)=\dfrac{2x-1}{x-3}$ \hfill \textbf{A1}

\medskip

\textbf{(e)} Domain of $f^{-1}$: $x\in\mathbb{R}$, $x\ne 3$. \hfill \textbf{A1}

\emph{FT: award A1 for ``$x\ne$ (the horizontal asymptote value of $f$)''
read from the candidate's answer in (b).}

\bigskip\hrule\bigskip

%================================================================
%--- Problem 2 [10 marks]
%================================================================
\noindent\textbf{2.} $\displaystyle g(x)=\dfrac{2-2x}{x+2}$, \ $x\ne -2$.

\medskip

\textbf{(a)} Let $y=g(x)$ and swap $x$ and $y$: \ \ $x=\dfrac{2-2y}{y+2}$
\hfill \textbf{M1}

$x(y+2)=2-2y$, \ \ \ $xy+2y=2-2x$, \ \ \ $y(x+2)=2-2x$ \hfill \textbf{A1}

$g^{-1}(x)=\dfrac{2-2x}{x+2}=g(x)$ \hfill \textbf{A1 AG}

\medskip

\textbf{(b)} \emph{Sketch via GDC (smooth two-branch reciprocal-type curve,
upper branch to the left of $x=-2$ heading to $y=-2$, lower branch to the right
heading to $y=-2$).}

Vertical asymptote: $x=-2$ \hfill \textbf{A1}

Horizontal asymptote: $y=-2$
\ \ (ratio of leading coefficients $=\tfrac{-2}{1}$) \hfill \textbf{A1}

$x$-intercept: $2-2x=0 \Rightarrow x=1$, \ so $(1,0)$ \hfill \textbf{A1}

$y$-intercept: $g(0)=\dfrac{2}{2}=1$, \ so $(0,1)$ \hfill \textbf{A1}

\emph{Accept $x=1$ and $y=1$ in place of the ordered pairs.}

\medskip

\textbf{(c)} For $h(x)=\dfrac{2-3x}{cx+d}$:

Vertical asymptote: $x=-\dfrac{d}{c}$ \hfill \textbf{A1}

Horizontal asymptote: $y=-\dfrac{3}{c}$ \hfill \textbf{A1}

\medskip

\textbf{(d)} \ A M\"obius function $\dfrac{ax+b}{cx+d}$ satisfies
$h^{-1}=h$ iff $a+d=0$. \\
With $a=-3$, this gives $d=3$. \hfill \textbf{A1}

\emph{Accept derivation by swapping $x\leftrightarrow y$ and requiring the
resulting expression to equal $h(x)$, leading to $d=-a=3$.}

\bigskip\hrule\bigskip

%================================================================
%--- Problem 3 [8 marks]
%================================================================
\noindent\textbf{3.} Transform $y=\dfrac{1}{x}$: \ vertical stretch $\times 2$ $\rightarrow$
reflection in $x$-axis $\rightarrow$ translation $4$ right and $3$ up. Image $y=p(x)$.

\medskip

\textbf{(a)} After the three transformations: \hfill \textbf{M1}

$\displaystyle y=\dfrac{2}{x}\ \to\ y=-\dfrac{2}{x}\ \to\ y=-\dfrac{2}{x-4}+3$

Combine over a common denominator: \hfill \textbf{M1}

$\displaystyle p(x)=\dfrac{-2+3(x-4)}{x-4}=\dfrac{3x-14}{x-4}$ \hfill \textbf{A1}
\ \ \ ($a=3,\ b=-14$).

\medskip

\textbf{(b)} Asymptotes: $x=4$ and $y=3$. \hfill \textbf{A1}

\emph{Both required for the A1.}

\medskip

\textbf{(c)} $x$-intercept: $3x-14=0 \Rightarrow x=\dfrac{14}{3}$ \hfill \textbf{A1}

$y$-intercept: $p(0)=\dfrac{-14}{-4}=\dfrac{7}{2}$ \hfill \textbf{A1}

\emph{Accept $\left(\tfrac{14}{3},0\right)$ and $\left(0,\tfrac{7}{2}\right)$.}

\medskip

\textbf{(d)} Let $y=\dfrac{3x-14}{x-4}$ and swap $x$ and $y$:

$x(y-4)=3y-14$, \ \ \ $y(x-3)=4x-14$ \hfill \textbf{M1}

$\displaystyle p^{-1}(x)=\dfrac{4x-14}{x-3}$, \ domain $x\in\mathbb{R}$, $x\ne 3$.
\hfill \textbf{A1}

\emph{FT: award A1 for $x\ne$ (the candidate's horizontal-asymptote value from (b)).}

\bigskip\hrule\bigskip

%================================================================
%--- Problem 4 [6 marks]
%================================================================
\noindent\textbf{4.} $f(x)=\dfrac{1}{x-1}$, \ $g(x)=2x+3$.

\medskip

\textbf{(a)} $(f\circ g)(x)=f(2x+3)=\dfrac{1}{(2x+3)-1}$ \hfill \textbf{M1}

$\displaystyle =\dfrac{1}{2x+2}=\dfrac{1}{2(x+1)}$ \hfill \textbf{A1}
\ \ ($a=2,\ h=-1$).

Domain: $x\in\mathbb{R}$, $x\ne -1$. \hfill \textbf{A1}

\medskip

\textbf{(b)} Re-express as $\displaystyle (f\circ g)(x)=\dfrac{1}{2}\cdot\dfrac{1}{x+1}$
\hfill \textbf{M1}

\emph{Two transformations of $y=\dfrac{1}{x}$:}

(i) horizontal translation of $1$ unit to the left
\ \ \ ($x\mapsto x+1$), \hfill \textbf{A1}

(ii) vertical stretch with scale factor $\dfrac{1}{2}$
\ \ \ (equivalently, a vertical compression with scale factor $2$). \hfill \textbf{A1}

\emph{Accept the two transformations applied in either order
(the order does not matter when one is horizontal and the other vertical).}

\end{document}
